% Local proof fragment.  Include in a document with amsmath, amsthm and url.
\begin{proposition}\label{prop:decic-no81}
An elliptic curve over a number field of degree $10$ has no point of
order $81$.
\end{proposition}
\begin{proof}
Put $X=X_1(81)$. We use the rank-zero result
\cite[Theorem~3.1]{DEHMZ} and the retained integral homology computation
\[
 J_X(\mathbb Q)[2]=0.
\]
For completeness, the latter computation uses the integral homology of
rank $380$. The kernel of $T_5-1-5\langle5\rangle$ modulo $2$ has
dimension $8$, and its intersection with the kernel of
$T_7-1-7\langle7\rangle$ is zero. Both operators kill rational
$2$-torsion. The retained matrices have been independently replayed in
\path{parent_two_torsion_replay.json}.

Reduction at $2$ is injective on $J_X(\mathbb Q)$. Indeed, the kernel
on torsion in $J_X(\mathbb Q_2)$ is killed by $2$: multiplication by $2$
maps the first formal filtration into the second, and the formal
logarithm is injective on the latter. The assertion follows from the
absence of rational $2$-torsion.

Suppose that a point of order $81$ is defined over a field of degree
$10$, and let $D$ be its trace divisor on $X$. Thus $D$ is rational,
effective, entirely noncuspidal and of degree $10$.

First suppose that $\overline D$ is cuspidal. The cusps are finite
\'{e}tale away from $81$, and their residue fields are subfields of
$\mathbb Q(\zeta_{81})$. Since $2$ generates
$(\mathbb Z/81\mathbb Z)^\times$, every global cusp orbit remains
intact on reduction at $2$. Hence $\overline D$ lifts to a rational
effective cuspidal divisor $C$ of degree $10$. Injectivity gives
$D\sim C$. Abramovich's bound~\cite{Abramovich} gives
\[
 \operatorname{gon}_{\mathbb C}X
 \geq \frac7{800}[\operatorname{PSL}_2(\mathbb Z):\overline{\Gamma_1(81)}]
 =\frac7{800}\,2916>25.
\]
Since $D\ne C$, the equivalence $D\sim C$ would give a function of
degree at most $10$, a contradiction.

We now consider reductions with noncuspidal support. The exhaustive
finite-field census gives the following numbers of noncuspidal closed
points:
\[
\begin{array}{c|rrrrrrrrrr}
 \text{degree}&1&2&3&4&5&6&7&8&9&10\\ \hline
 \text{number}&0&0&0&0&0&0&0&0&135&135.
\end{array}
\]
In each of degrees $9$ and $10$ these points form five orbits under
$(\mathbb Z/81\mathbb Z)^\times/\{\pm1\}$. Thus, after a rational
diamond operator, $\overline D$ is one of five degree-$10$ points, or
one of five degree-$9$ points together with a rational cusp. We number
these ten noncuspidal supports as in \path{source81.m}, with IDs
$1,\ldots,5$ in degree $9$ and $6,\ldots,10$ in degree $10$.

The census enumerates every elliptic curve isomorphism class over
$\mathbb F_{2^d}$ for $d\leq10$, and every cyclic subgroup of order
$81$ on the relevant curves. In degree $6$ the only group of order
divisible by $81$ is $C_9\times C_9$, which contributes no point of
order $81$. The degree-$9$ contribution consists of $45$ curves with
group $C_{486}$. In degree $10$ there are $30$ relevant curves, with
groups $C_{972}$ or $C_3\times C_{324}$, and $50$ cyclic subgroups of
order $81$. The second enumeration in
\path{parent_independent_source_filter.json} independently recovers
all ten normalized records and both numbers $135$.

The map $X\longrightarrow X_0(81)$ has odd degree $27$. Consequently
$J_0(81)(\mathbb Q)$ has rank zero and has no rational $2$-torsion.
The proven finite Jacobian presentation is
\[
 J_0(81)(\mathbb F_2)\simeq C_3\times C_{18}.
\]
Its odd subgroup has order $27$ and contains the reduction of every
rational divisor class of degree zero. The normalized noncuspidal
classes for IDs $4,5,6,8,9$ lie outside this subgroup. A rational
cuspidal complement has class in the subgroup and cannot change this
obstruction. Every inverse image of the computed pair of $j$-invariants
was included, with its norm multiplicity; all ten inverse images are
unique and have norm multiplicity $1$. This leaves precisely
\[
                  1,2,3,7,10.                         \tag{81.1}
\]
The finite and infinite ideal identities for this sieve were replayed
on a fresh native model in \path{parent_native_x0_filter.m}.

It remains to exclude (81.1). Put
\[
 Y=X_1(81)/\langle\langle28\rangle\rangle,
 \qquad Z=X_1(27),
\]
so that $X\longrightarrow Y\longrightarrow Z$ has degrees $3,3$.
Here $g(Z)=13$ and $g(Y)=46$. The rank-zero assertion and the absence
of rational $2$-torsion descend from $J_X$ to $J_Y$: pullback followed
by norm is multiplication by the odd integer $3$. The curve $Y$ has
good reduction at $2$, since it is the tame quotient of the smooth
model of $X$ by a group of order $3$. Reduction at $2$ is therefore
injective on $J_Y(\mathbb Q)$.

Write $\delta=\langle2\rangle$. The Eichler--Shimura relation gives
\[
             (T_2-1-2\delta)J_Y(\mathbb Q)=0.           \tag{81.2}
\]
For the rational Tate cusp $c_0$ we have the literal identity
$T_2c_0=c_0+2\delta c_0$. Indeed, quotienting
$(\operatorname{Tate}(q),q^{1/81})$ by $\mu_2$ preserves its component
label; the other two subgroups give the component label $2/81$.
The same identity holds after passing to $Y$.

Let $E$ be the pushforward of $D$ to $Y$. Applying (81.2) to
$[E-10c_0]$ gives
\[
                   T_2E\sim E+2\delta E.                \tag{81.3}
\]
Both divisors have degree $30$. On the special fibre, the Hecke
correspondence acts on cycles by Frobenius plus $\delta$ times
Frobenius pullback. Since $\overline E$ is rational over $\mathbb F_2$,
\[
              \overline{T_2E}=\overline E+2\delta\overline E.
                                                               \tag{81.4}
\]
We compute, for every remaining possibility,
\[
              h^0_Y(\overline E+2\delta\overline E)=1.   \tag{81.5}
\]
We give the model and the finite linear algebra for this computation
below. Semicontinuity now gives $h^0_Y(E+2\delta E)=1$. Hence (81.3)
is an equality of effective divisors.

The finite set of $j$-invariants in the support of $E$ is therefore
closed under every cyclic $2$-isogeny. This is impossible. For
$E_\tau=\mathbb C/(\mathbb Z+\tau\mathbb Z)$, the map $z\mapsto2z$
gives a cyclic $2$-isogeny $E_\tau\longrightarrow E_{2\tau}$.
Thus all $j(2^k\tau)$
would lie in the same finite set, whereas these numbers are unbounded.
This proves the proposition.
\end{proof}

\paragraph{The genus-$46$ computation.}
We describe the calculation of (81.5), including the constant field,
the selected points in the cubic cover and the conditions at the
other branches. All function-field data are saved in
\path{y46_kummer_model.m}.

Let $(\mathcal E,R)$ be the universal elliptic curve with a point of
order $27$, in Tate normal form. Put $S=9R$, let $m$ be the tangent
slope at $S$, and set
\[
                f=(y(R)-y(S))-m(x(R)-x(S)).
\]
Over $\mathbb F_4$, the cover $Y\longrightarrow Z$ is
\[
                         w^3=f.                         \tag{81.6}
\]
To identify the twist in (81.6), take
$\phi:\mathcal E\longrightarrow\mathcal E/\langle S\rangle$ and
its dual $\widehat\phi$. The curve $Y$ parametrizes points $V$ with
$\widehat\phi(V)=R$. This follows directly from
$(\mathcal E,P)\mapsto(\mathcal E,3P,\phi(P))$; replacing $P$ by
$28P$ does not change this image. The cubic equation for
$x(V)$ is irreducible over $\mathbb F_4(Z)$, and the retained
functions $a,A,B$ satisfy the exact identity
\[
                         x(V)=a+Aw+Bw^2.
\]
Substitution in the cubic verifies (81.6) and its twist. The divisor
of $f$ has eleven rational places over $\mathbb F_4$, with
multiplicities
\[
             2,2,8,-9,-22,-9,-7,14,8,8,5.
\]
Exactly nine places are branched. All are cusps and have valuation
$2$ modulo $3$. Riemann--Hurwitz gives
$g(Y)=1+3(13-1)+9=46$.

The nine rational source-cusp images are these nine ramified cusps.
Their images in $Z$ are rational over $\mathbb F_2$ and are permuted
in one $9$-cycle by $\delta$. This is also determined by cusp widths:
the map to $Z$ has ramification index $3$ at the images of the
rational Tate cusps. The computed permutation, in the retained order,
is
\[
                  [2,8,6,1,9,4,5,7,3].
\]
Thus, for each of IDs $1,2,3$, we test all nine choices of the
cuspidal complement. IDs $7,10$ have no complement. Base change to
$\mathbb F_4$ preserves $h^0$; a degree-$9$ noncuspidal point stays
irreducible, while a degree-$10$ point splits into two degree-$5$
points. All these points have norm multiplicity $1$ in the successive
maps used here.

Write $\overline E=N+C$ and put
$L=N+2\delta N+C+2\delta C$ on $Y_{\mathbb F_4}$. Let $M$ be the
pushforward of $N+2\delta N$ to $Z$. Its support is disjoint from the
branch locus. First allow poles of order $M_p$ on all three points
above every $p\in\operatorname{supp}(M)$, and retain the prescribed
poles at the ramified cusps. The space with these relaxed conditions
decomposes into the three characters of the cover. Its character-$j$
part is
\[
 H^0\left(Z,M+
       \sum_p\left\lfloor\frac{c_p+jv_p(f)}3\right\rfloor p\right)w^j,
       \qquad j=0,1,2,                                  \tag{81.7}
\]
where $c_p$ is the coefficient of $C+2\delta C$ at the unique
ramified point above $p$, and is zero otherwise. The sum includes
all zeros and poles of $f$.

At a noncuspidal place $p$ the three cube roots of $f(p)$ are
distinct. We select the actual root by the displayed formula for
$x(V)$ and a fresh $3$-isogeny calculation on the source elliptic
curve. If $\tau$ is a uniformizer at $p$ and the allowed pole order
is $r\in\{1,2\}$, we impose vanishing to order $r$ of
$\tau^r\sum_jg_jw^j$ at each of the other two roots. Constant terms
give the conditions for $r=1$; constant and first terms give the
conditions for $r=2$. These are all the omitted local conditions, so
their kernel is exactly $H^0(Y,L)$.

The dimensions and matrix ranks are
\[
\begin{array}{c|c|c|c|c|c}
 \text{IDs}&\text{tests}&\text{character dimensions}
       &\text{matrix size}&\text{rank}&h^0\\ \hline
 1,2,3&27&(15,11,13)&39\times54&38&1\\
 7,10&2 &(18,12,15)&45\times60&44&1.
\end{array}
\]
This proves (81.5). The files \path{rr_case1_summary.m},
\path{rr_case2_summary.m}, \path{rr_case3_summary.m},
\path{rr_case7_summary.m} and \path{rr_case10_summary.m} record all
$29$ tests.

The native replay \path{parent_native_completion_RR_v2.m} rebuilds
all $87$ spaces in (81.7), checks membership and independence of
every exported basis, and recomputes every matrix entry from genuine
local completions. It also checks the uniformizers and selects the
physical branch from the source level point. Independently, the Sage
replay in \path{independent_sage_local_arcs.json} constructs Hensel
arcs to precision $64$, checks the source point orders and the
explicit V\'{e}lu images, and recovers every entry of all $29$ matrices.
Both replays give the ranks in the table.
